\documentclass[12pt]{article}
\usepackage{amsmath, amssymb, geometry, setspace}
\geometry{margin=1in}
\setstretch{1.2}
\setlength{\parindent}{0pt}


\title{Quiz 3\\ \vspace{10pt}
\large ECON 320 -- Introduction to Mathematical Economics}
\author{Instructor: Muhebullah Karimzada}
\date{October 20, 2025}

\begin{document}
\maketitle



\section*{Part A – Multiple Choice Questions (1 point each)}

\begin{enumerate}
\item If \(f(x)\) has a local maximum at \(x=a\), then  
\begin{enumerate}
\item \(f'(a)>0\)  
\item \(f'(a)=0\) and \(f''(a)>0\)  
\item \(f'(a)=0\) and \(f''(a)<0\)  
\item \(f'(a)<0\)
\end{enumerate}

\item The second-order condition for a \textit{minimum} of a differentiable function \(f(x)\) is  
\begin{enumerate}
\item \(f'(x)=0\) and \(f''(x)<0\)  
\item \(f'(x)=0\) and \(f''(x)>0\)  
\item \(f''(x)=0\)  
\item \(f'(x)>0\)
\end{enumerate}

\item If \(f(x)=x^3-6x^2+9x\), the function attains a minimum at  
\begin{enumerate}
\item \(x=1\)  
\item \(x=2\)  
\item \(x=3\)  
\item \(x=0\)
\end{enumerate}

\newpage

\item The profit function \(\pi(Q)=120Q-5Q^2\) reaches its maximum when marginal revenue equals  
\begin{enumerate}
\item 120  
\item zero  
\item 5  
\item average cost
\end{enumerate}

\item The function \(y=x\,e^{-x}\) has its maximum value when  
\begin{enumerate}
\item \(x=0\)  
\item \(x=1\)  
\item \(x=2\)  
\item \(x=1/e\)
\end{enumerate}

\item If \(f(x)=a-bx^2\) with \(a,b>0\), then \(f(x)\) is  
\begin{enumerate}
\item strictly concave  
\item strictly convex  
\item linear  
\item neither concave nor convex
\end{enumerate}

\item The necessary condition for an interior extremum of \(f(x)\) is that  
\begin{enumerate}
\item \(f'(x)\neq0\)  
\item \(f''(x)=0\)  
\item \(f'(x)=0\)  
\item \(f''(x)>0\)
\end{enumerate}
\end{enumerate}

\newpage

\section*{Part B – Analytical Problems (13 marks total)}

\textbf{Question 1 (4 points)}\\
A firm’s profit is given by  
\[
\pi(Q)=100Q-4Q^2-20.
\]
\begin{enumerate}
\item Derive the first-order condition.  
\item Find the profit-maximizing output \(Q^*\).  
\item Compute the maximum profit.  
\item Verify the second-order condition.
\end{enumerate}

\vspace{1em}

\textbf{Question 2 (4 points)}\\
Let the utility of consumption be  
\[
U(x)=10x-0.5x^2.
\]
\begin{enumerate}
\item Find \(x^*\) that maximizes utility.  
\item Check whether this is a maximum.  
\item Evaluate \(U(x^*)\).  
\item Explain the economic meaning of the result.
\end{enumerate}

\vspace{1em}

\textbf{Question 3 (5 points)}\\
The total revenue of a monopolist is  
\[
R(P)=500P e^{-0.1P}.
\]
\begin{enumerate}
\item Derive the price \(P^*\) that maximizes total revenue.  
\item Compute the corresponding maximum revenue \(R(P^*)\).  
\item Verify that the second-order condition holds.  
\item Discuss why \(R(P)\) first rises and then falls as \(P\) increases.  
\end{enumerate}

\end{document}
